\section{Benchmark design}

\subsection{Task definition}

Given source MIDI $x$, instruction $e$, a system produces edited MIDI $\hat{x}$ and optionally a structured plan $\hat{p}$. The instruction specifies attributes that \emph{must change}; all other attributes \emph{should preserve} unless the instruction explicitly overrides them.

Each benchmark item stores:
\begin{itemize}[leftmargin=*]
\item paths to $x$ and (for unique gold) reference $x^\*$;
\item instruction $e$ and operator family label;
\item edit mask $m$ (tracks and optional time window);
\item lists of \texttt{must\_change} and \texttt{must\_preserve} predicates;
\item optional gold plan $p$ for MIDI-Reason.
\end{itemize}

\subsection{Gold modes}

\paragraph{Unique gold (v0.1).}
Apply deterministic transform $T_\theta$ to seed MIDI, verbalize $\theta$ as $e$, store $x^\* = T_\theta(x)$. Enables note-F1 and exact plan matching.

\paragraph{Constraint gold (planned).}
For descriptive instructions with many valid outputs, store necessary unit tests (Not that Groove style) without a single $x^\*$. Deferred until unique-gold grader validation completes.

\subsection{MIDI-Reason track}

The plan $p = (\text{op}, \theta)$ is scored with exact match independently from $\hat{x}$. This separates \emph{naming} the edit from \emph{executing} it---the same distinction agent security benchmarks draw between belief probes and tool-call behavior.

\subsection{Operator families (v0.1)}

Transpose (non-drum tracks), tempo scale, velocity scale, track mute, and program change. All operators are invertible or programmatically verifiable; no LLM judge is used for headline metrics.
